\subsection{Longer training budget for the serial credit comparison}

The original physical-depth comparisons used a maximum of 180 epochs and validation-loss early stopping with patience thirty. An audit of one hundred original D1/D3 trajectories verified their configurations and final metrics against the archived files. All fifty D3 fits across the five relevant credit/optimization conditions reached 180 epochs, with validation loss still decreasing. These endpoints therefore cannot establish a converged advantage.

The follow-up retained all ten original seeds in five aligned D3 conditions: exact autograd and soma-broadcast autograd under the BP recipe, soma-broadcast autograd under the LocalCA recipe, and exact-path and shared-soma LocalCA. Ten exact-BP D1 fits supplied the depth reference, giving sixty fits. Every model had 66,178 trainable parameters and 14,336 active synapses. Both recipes used Adam; the BP recipe used learning rate 0.003 without split parameter groups, while the LocalCA recipe retained its original split groups and nominal rate 0.0008. Complete group-specific rates, weight-boosting, decay, clipping and update options are retained in the configurations. Only the maximum epoch count and output names changed. The original patience and validation-based state selection were preserved, without tuning against the extended outcomes.

Historical final-model files did not contain optimizer and random-generator states. The extension therefore restarted the original seeds using the pinned implementation \texttt{a99c3a7}, rather than claiming exact continuation from epoch 180. Four seed blocks used GPUs and six used CPUs, with the device class held fixed across all six conditions within a seed. At epochs 180, 360 and 600, and on earlier termination, additional records retained current and best weights, optimizer state and Python, NumPy and Torch random states. A separate audit confirmed that every final model exactly equaled its saved validation-selected best state, and that all configuration changes were limited to the epoch cap and output naming.

For each restart, we evaluated the best validation-selected state available within the first 180 epochs and within the complete run. If a run stopped before 180, both windows used its final selected state. Test accuracy and cross-entropy were reporting metrics, never selection criteria. Learning curves show best validation loss; stopped runs are carried forward so every displayed mean includes the same ten seeds. The compact Source Data retain all 34,159 observed epochs, every recorded epoch metric, 120 original/extension configurations, sixty run records and the complete underlying raw-file hash inventory. Carried plotting rows are separately identified.

At 600 epochs, the exact-BP D3-minus-D1 accuracy difference was 38.17 percentage points (95\% paired interval, 37.98--38.39). The LocalCA exact-path-minus-shared-soma difference changed from 10.86 points (10.15--11.62) at 180 epochs to $-1.52$ ($-1.83$ to $-1.23$) at 600. The BP-recipe exact-minus-broadcast difference at 600 was 0.48 points (0.40--0.55), and the broadcast BP-minus-LocalCA-recipe difference was 1.54 points (1.39--1.70). Secondary loss contrasts retain the metric dependence: exact-path LocalCA had mean test cross-entropy 0.2544 versus 0.2758 for shared-soma LocalCA, despite lower accuracy. All contrasts use the same ten paired seeds and 10,000 whole-seed bootstrap draws; intervals are descriptive, without a new confirmatory multiplicity claim.

All fifty D3 fits and two D1 fits reached the 600-epoch cap. Every D3 condition still had decreasing validation loss over its final thirty epochs. Relative to the historical first-180 trajectories, the largest absolute validation-loss difference was 0.005885. The median absolute discrepancy in validation-selected test accuracy was 0.070 percentage points and the largest was 1.830 points. These differences are retained per run; budget contrasts use only the two windows of the same new trajectories. The extension establishes budget and metric dependence, not convergence or fresh-seed replication.
